\section{Forces and the Higgs} \label{sec:forces}

The same division-algebra tower that capped the dimension at eight also names the three forces. Each rung contributes
one factor of the gauge group, and the Higgs is the connection along the quaternionic direction. Every dimension count
below is checked by computation.

\subsection{The gauge group, rung by rung}

A continuous symmetry that preserves a normed algebra is its group of norm-preserving automorphisms. Walk the tower
and read these groups off.

\paragraph{The complexes give $U(1)$.} The unit-norm elements of $\mathbb{C}$ are the unit circle, the phases
$e^{i\theta}$, which form the group $U(1)$. This is the hypercharge factor, the abelian phase rotation, dimension one.

\paragraph{The quaternions give $SU(2)$.} The unit-norm quaternions are the three-sphere $S^3$, and as a group they
are exactly $SU(2)$, the double cover of rotations, dimension three. This is the weak isospin factor. We verify the
identification by exhibiting the standard isomorphism on the basis.

\paragraph{The octonions give $SU(3)$.} The automorphism group of the octonions is $G_2$, dimension fourteen. Fixing
one imaginary unit, equivalently choosing a complex direction inside the octonions and asking which automorphisms
preserve it, leaves the subgroup that fixes that unit, and the stabilizer of an imaginary unit in $G_2$ is $SU(3)$, of
dimension eight. We verify the arithmetic, $G_2$ has dimension fourteen, the orbit of a unit imaginary is the
six-sphere of dimension six, and $14 - 6 = 8 = \dim SU(3)$. This is the color factor.

Collecting the three,
\begin{equation}
  SU(3) \times SU(2) \times U(1), \qquad \dim = 8 + 3 + 1 = 12 .
\end{equation}

\begin{result}[The gauge group]
The Standard Model gauge group $SU(3) \times SU(2) \times U(1)$ is the product of the symmetry groups of the three
non-trivial composition algebras, the octonions giving color $SU(3)$ as the unit-fixing subgroup of $G_2$, the
quaternions giving weak $SU(2)$ as the unit three-sphere, and the complexes giving hypercharge $U(1)$ as the unit
circle. Its total dimension is twelve.
\end{result}

The acting of this group on the fermion Fock space of Section~\ref{sec:fermions} is the standard one. The $SU(3)$
rotates the three color modes, the $SU(2)$ acts on the complex-doublet structure $\mathbb{C} \otimes \mathbb{O}$
restricted to the quaternionic part, and the $U(1)$ is the overall phase, with hypercharge fixed by the number
operator. The chirality, that the weak force acts only on left-handed states, comes from the two half-spinors $8_s$ and
$8_c$ being inequivalent, the left and right of the world being genuinely different representations, which is again a
fact special to triality.

\subsection{The Higgs is the quaternionic doublet}

In the noncommutative-geometry reading of the Standard Model \cite{connes1996}, the Higgs field is not an extra
ingredient but the part of the connection that points in the internal, discrete direction, exactly as the gauge fields
are the parts pointing in spacetime directions. The internal direction here is the quaternionic one, and the
quaternions are isomorphic as a complex vector space to $\mathbb{C}^2$, a two-component complex object, a doublet of
$SU(2)$.

\begin{result}[The Higgs]
The Higgs field is the connection along the quaternionic internal direction, which as an $SU(2)$ representation is the
doublet $\mathbf{2}$, since $\mathbb{H} \cong \mathbb{C}^2$. It is a doublet and not a triplet, which is required by
the measured custodial relation. A doublet gives the tree-level ratio $\rho = M_W^2 / (M_Z^2 \cos^2\theta_W) = 1$,
while a triplet would give $\rho = \tfrac12$.
\end{result}

We verify the custodial arithmetic. For a scalar in an isospin-$T$, hypercharge-$Y$ representation acquiring a vacuum
value, the tree-level $\rho$ parameter is
\begin{equation}
  \rho = \frac{T(T+1) - Y^2}{2 Y^2},
\end{equation}
which equals one exactly for the doublet $T = \tfrac12, Y = \tfrac12$, and equals one half for the triplet $T = 1$.
Experiment gives $\rho = 1$ to a fraction of a percent, so the doublet is selected. The single scalar doublet of the
Standard Model is therefore the quaternionic connection, its two complex components being the two halves of
$\mathbb{H} \cong \mathbb{C}^2$.

So the three rungs above the integers, the complexes, quaternions, and octonions, give the three forces, and the
quaternionic rung gives, in addition, the one scalar that breaks the symmetry. The forces, the matter, and the Higgs
all issue from the single tower that the pinch-point capped at eight.
